%% refer q31.tex
%% include the topic name as per u r project in appendix

\appendix
\chapter{API Reference}

This appendix provides a complete reference for all REST API endpoints exposed by the Smart Queue Management System backend. All endpoints are prefixed with \texttt{http://localhost:5000}. Request and response bodies are in JSON format.

\begin{table}[htbp]
    \centering
    \begin{tabular}{|c|c|c|c|}
        \hline
        \textbf{Method} & \textbf{Endpoint} & \textbf{Request Body} & \textbf{Description} \\
        \hline
        POST & /api/token/generate & \{customerName\} & Generate a new queue token \\
        \hline
        GET & /api/token/all?startDate=YYYY-MM-DD\&endDate=YYYY-MM-DD & None & Get all tokens or filter by date range \\
        \hline
        GET & /api/token/current & None & Get the active token \\
        \hline
        GET & /api/token/stats & None & Get waiting, active, completed count \\
        \hline
        POST & /api/token/next & None & Call the next waiting token \\
        \hline
        PUT & /api/token/complete/\{id\} & None & Mark a token as completed \\
        \hline
        DELETE & /api/token/reset & None & Delete all tokens \\
        \hline
        POST & /api/admin/login & \{username, password\} & Admin authentication \\
        \hline
    \end{tabular}
    \caption{Complete REST API Reference}
    \label{tab:apiref}
\end{table}

\chapter{Database Schema}

The following SQL script can be used to manually create the database tables for the Smart Queue Management System. Under normal operation, the tables are created automatically by Entity Framework Core when the backend application starts.

\begin{verbatim}
-- Create the database
CREATE DATABASE queue_management;

-- Connect to the database
\c queue_management

-- Create Admins table
CREATE TABLE IF NOT EXISTS "Admins" (
    "Id"       SERIAL PRIMARY KEY,
    "Username" VARCHAR(100) NOT NULL UNIQUE,
    "Password" VARCHAR(100) NOT NULL
);

-- Create Tokens table
CREATE TABLE IF NOT EXISTS "Tokens" (
    "Id"           SERIAL PRIMARY KEY,
    "TokenNumber"  INT          NOT NULL,
    "CustomerName" VARCHAR(200) NOT NULL,
    "Status"       VARCHAR(50)  NOT NULL DEFAULT 'Waiting',
    "CreatedAt"    TIMESTAMP    NOT NULL DEFAULT NOW()
);

-- Seed default admin user
INSERT INTO "Admins" ("Username", "Password")
VALUES ('admin', 'admin123')
ON CONFLICT ("Username") DO NOTHING;
\end{verbatim}
